% GENERATED FILE. Edit canonical lessons, then run npm run generate:books.
% canonical-source-hash: fnv1a64:72eacef19f3539dd
% canonical-lessons: AR-C302-hiwar, AR-C302-hawara, AR-C302-ajra, AR-C302-istadafa, AR-C302-mutahaddith

\chapter{An interview}
\label{ch:ar-an-interview}
\hlchaptermodality{302}

\begin{chapteropening}
\textbf{By the end of this chapter:} \emph{I can say a dialogue, to interview, to conduct, to host and a spokesperson.}

Complete the last lesson of chapter 302: I can say a dialogue, to interview, to conduct, to host and a spokesperson.
\end{chapteropening}

\section[\texorpdfstring{\ar{حوار} (ḥiwār)}{حوار (ḥiwār)}]{\ar{حوار} (ḥiwār) — a dialogue, a conversation}

\label{lesson:AR-C302-hiwar}



\begin{quote}
Before the new one: say the Arabic for a wounded person, then the Arabic for to give first aid.
\end{quote}

\subsection*{You'll want to know: \ar{حوار}}
\textbf{\ar{حوار}} — \emph{ḥiwār} — \textquotedblleft{}a dialogue, a conversation\textquotedblright{}.

\begin{quote}
\ar{نشرت} \ar{المجلة} \ar{حوارا} \ar{مع} \ar{الكاتب}.
\end{quote}

\begin{quote}
\emph{nasharat al-majalla ḥiwāran maʿa al-kātib.}
\end{quote}

\begin{quote}
The magazine published a conversation with the writer.
\end{quote}

In the press a \ar{حوار} is an interview printed as questions and answers.

\begin{cousinweb}
\textbf{\ar{حوار}} (\emph{ḥiwār}) is built on the native Arabic root \emph{ḥ-w-r} (Lane), \textquotedblleft{}to go back, to return\textquotedblright{}. Speech that goes back and forth is a \emph{ḥiwār}, the verbal noun of Form III.
\end{cousinweb}

\begin{grammarlens}[title={What you've built}]
One more word for telling what happened.
\end{grammarlens}

\subsection*{Your turn}
\begin{itemize}
\raggedright
  \item \emph{Say it:} \emph{ḥiwār}
  \item \emph{Say it:} \emph{ḥiwār}, once more
  \item \emph{Say it:} \emph{nasharat al-majalla ḥiwāran maʿa al-kātib.} — The magazine published a conversation with the writer.
  \item \emph{Recall:} say the Arabic for a wounded person, then the Arabic for to give first aid, then say \emph{ḥiwār} again
\end{itemize}

\begin{tcolorbox}[breakable,colback=teal!4,colframe=teal!35!black,title={Before you move on}]
What does \ar{حوار} mean? (\textquotedblleft{}A dialogue, a conversation\textquotedblright{}.) Say it once more.
\end{tcolorbox}
\section[\texorpdfstring{\ar{حاور} (ḥāwara)}{حاور (ḥāwara)}]{\ar{حاور} (ḥāwara) — to hold a dialogue; to interview}

\label{lesson:AR-C302-hawara}



\begin{quote}
Before the new one: say the Arabic for to give first aid, then the Arabic for a dialogue.
\end{quote}

\subsection*{You'll want to know: \ar{حاور}}
\textbf{\ar{حاور}} — \emph{ḥāwara} — \textquotedblleft{}to hold a dialogue; to interview\textquotedblright{}.

\begin{quote}
\ar{حاور} \ar{المذيع} \ar{لاعبا} \ar{مشهورا}.
\end{quote}

\begin{quote}
\emph{ḥāwara al-mudhīʿ lāʿiban mashhūran.}
\end{quote}

\begin{quote}
The presenter interviewed a famous player.
\end{quote}

Present: \ar{يحاور}, \emph{yuḥāwir}.

\begin{cousinweb}
\textbf{\ar{حاور}} (\emph{ḥāwara}) is built on the native Arabic root \emph{ḥ-w-r} (Lane), \textquotedblleft{}to go back\textquotedblright{}. It is Form III, \textquotedblleft{}to answer one another\textquotedblright{}.
\end{cousinweb}

\begin{grammarlens}[title={What you've built}]
One more word for telling what happened.
\end{grammarlens}

\subsection*{Your turn}
\begin{itemize}
\raggedright
  \item \emph{Say it:} \emph{ḥāwara}
  \item \emph{Say it:} \emph{ḥāwara}, once more
  \item \emph{Say it:} \emph{ḥāwara al-mudhīʿ lāʿiban mashhūran.} — The presenter interviewed a famous player.
  \item \emph{Recall:} say the Arabic for to give first aid, then the Arabic for a dialogue, then say \emph{ḥāwara} again
\end{itemize}

\begin{tcolorbox}[breakable,colback=teal!4,colframe=teal!35!black,title={Before you move on}]
What does \ar{حاور} mean? (\textquotedblleft{}To hold a dialogue; to interview\textquotedblright{}.) Say it once more.
\end{tcolorbox}
\section[\texorpdfstring{\ar{أجرى} (ajrā)}{أجرى (ajrā)}]{\ar{أجرى} (ajrā) — to conduct, to carry out}

\label{lesson:AR-C302-ajra}



\begin{quote}
Before the new one: say the Arabic for a dialogue, then the Arabic for to hold a dialogue.
\end{quote}

\subsection*{You'll want to know: \ar{أجرى}}
\textbf{\ar{أجرى}} — \emph{ajrā} — \textquotedblleft{}to conduct, to carry out\textquotedblright{}.

\begin{quote}
\ar{أجرت} \ar{الصحفية} \ar{مقابلة} \ar{مع} \ar{الطبيب}.
\end{quote}

\begin{quote}
\emph{ajrat aṣ-ṣaḥafiyya muqābala maʿa aṭ-ṭabīb.}
\end{quote}

\begin{quote}
The journalist did an interview with the doctor.
\end{quote}

Present: \ar{يجري}, \emph{yujrī}.

\ar{أجرى} \ar{مقابلة}, to conduct an interview; \ar{أجرى} \ar{عملية}, to perform an operation.

\begin{cousinweb}
\textbf{\ar{أجرى}} (\emph{ajrā}) is built on the native Arabic root \emph{j-r-y} (Lane), \textquotedblleft{}to run, to flow\textquotedblright{}, the root of \emph{jarā}, to run. It is Form IV, \textquotedblleft{}to make a thing run\textquotedblright{}: to carry it through.
\end{cousinweb}

\begin{grammarlens}[title={What you've built}]
One more word for telling what happened.
\end{grammarlens}

\subsection*{Your turn}
\begin{itemize}
\raggedright
  \item \emph{Say it:} \emph{ajrā}
  \item \emph{Say it:} \emph{ajrā}, once more
  \item \emph{Say it:} \emph{ajrat aṣ-ṣaḥafiyya muqābala maʿa aṭ-ṭabīb.} — The journalist did an interview with the doctor.
  \item \emph{Recall:} say the Arabic for a dialogue, then the Arabic for to hold a dialogue, then say \emph{ajrā} again
\end{itemize}

\begin{tcolorbox}[breakable,colback=teal!4,colframe=teal!35!black,title={Before you move on}]
What does \ar{أجرى} mean? (\textquotedblleft{}To conduct, to carry out\textquotedblright{}.) Say it once more.
\end{tcolorbox}
\section[\texorpdfstring{\ar{استضاف} (istaḍāfa)}{استضاف (istaḍāfa)}]{\ar{استضاف} (istaḍāfa) — to host, to have as a guest}

\label{lesson:AR-C302-istadafa}



\begin{quote}
Before the new one: say the Arabic for to hold a dialogue, then the Arabic for to conduct.
\end{quote}

\subsection*{You'll want to know: \ar{استضاف}}
\textbf{\ar{استضاف}} — \emph{istaḍāfa} — \textquotedblleft{}to host, to have as a guest\textquotedblright{}.

\begin{quote}
\ar{استضافت} \ar{المدرسة} \ar{كاتبا} \ar{معروفا}.
\end{quote}

\begin{quote}
\emph{istaḍāfat al-madrasa kātiban maʿrūfan.}
\end{quote}

\begin{quote}
The school hosted a well-known writer.
\end{quote}

Present: \ar{يستضيف}, \emph{yastaḍīf}.

\begin{cousinweb}
\textbf{\ar{استضاف}} (\emph{istaḍāfa}) is built on the native Arabic root \emph{ḍ-y-f} (Lane), the root of the taught \emph{ḍayf}, a guest. It is Form X, first \textquotedblleft{}to ask someone for protection or refuge\textquotedblright{}; today a programme or a city that hosts.
\end{cousinweb}

\begin{grammarlens}[title={What you've built}]
One more word for telling what happened.
\end{grammarlens}

\subsection*{Your turn}
\begin{itemize}
\raggedright
  \item \emph{Say it:} \emph{istaḍāfa}
  \item \emph{Say it:} \emph{istaḍāfa}, once more
  \item \emph{Say it:} \emph{istaḍāfat al-madrasa kātiban maʿrūfan.} — The school hosted a well-known writer.
  \item \emph{Recall:} say the Arabic for to hold a dialogue, then the Arabic for to conduct, then say \emph{istaḍāfa} again
\end{itemize}

\begin{tcolorbox}[breakable,colback=teal!4,colframe=teal!35!black,title={Before you move on}]
What does \ar{استضاف} mean? (\textquotedblleft{}To host, to have as a guest\textquotedblright{}.) Say it once more.
\end{tcolorbox}
\section[\texorpdfstring{\ar{متحدث} (mutaḥaddith)}{متحدث (mutaḥaddith)}]{\ar{متحدث} (mutaḥaddith) — a spokesperson; a speaker}

\label{lesson:AR-C302-mutahaddith}



\begin{quote}
Before the new one: say the Arabic for to conduct, then the Arabic for to host.
\end{quote}

\subsection*{You'll want to know: \ar{متحدث}}
\textbf{\ar{متحدث}} — \emph{mutaḥaddith} — \textquotedblleft{}a spokesperson; a speaker\textquotedblright{}.

\begin{quote}
\ar{قال} \ar{متحدث} \ar{إن} \ar{المصنع} \ar{مغلق} \ar{اليوم}.
\end{quote}

\begin{quote}
\emph{qāla mutaḥaddith inna al-maṣnaʿ mughlaq al-yawm.}
\end{quote}

\begin{quote}
A spokesman said that the factory is closed today.
\end{quote}

\ar{متحدث} \ar{باسم} \ar{الشركة}: a spokesman for the company.

\begin{cousinweb}
\textbf{\ar{متحدث}} (\emph{mutaḥaddith}) is built on the native Arabic root \emph{ḥ-d-th} (Lane), the root of the taught \emph{ḥadīth}, talk. \emph{Mutaḥaddith} is the Form V participle, \textquotedblleft{}one who speaks\textquotedblright{}.
\end{cousinweb}

\begin{grammarlens}[title={What you've built}]
One more word for telling what happened.
\end{grammarlens}

\subsection*{Your turn}
\begin{itemize}
\raggedright
  \item \emph{Say it:} \emph{mutaḥaddith}
  \item \emph{Say it:} \emph{mutaḥaddith}, once more
  \item \emph{Say it:} \emph{qāla mutaḥaddith inna al-maṣnaʿ mughlaq al-yawm.} — A spokesman said that the factory is closed today.
  \item \emph{Recall:} say the Arabic for to conduct, then the Arabic for to host, then say \emph{mutaḥaddith} again
\end{itemize}

\begin{tcolorbox}[breakable,colback=teal!4,colframe=teal!35!black,title={Before you move on}]
What does \ar{متحدث} mean? (\textquotedblleft{}A spokesperson; a speaker\textquotedblright{}.) Say it once more.
\end{tcolorbox}
